\subsection{Matriz comparativa del escenario 3}

\begin{center}
\footnotesize
\begin{tabularx}{\textwidth}{|>{\raggedright\arraybackslash}p{2.6cm}|L|L|L|}
\hline
\rowcolor{verdeuteq}
\textcolor{white}{\textbf{Criterio}} &
\textcolor{white}{\textbf{Secuencial (cascada)}} &
\textcolor{white}{\textbf{Iterativo e incremental}} &
\textcolor{white}{\textbf{Ágil}} \\ \hline
K1. Estabilidad y cambios &
La ERS debe aprobarse en el mes 2. Los cambios del flujo de patentes y la posible ordenanza obligarían a tramitar modificaciones formales después de esa aprobación. &
Permite revisar en ciclos el flujo de patentes y adaptar el cálculo de tasas si cambia la ordenanza, sin perder el control de la versión aprobada. &
Ayudaría a ajustar el flujo y las tasas con rapidez, pero los cambios de alcance siguen sujetos a trámites contractuales que pueden durar semanas. \\ \hline
K2. Partes interesadas y disponibilidad &
Las aprobaciones por fase se pueden coordinar con el administrador del contrato, que es el único canal oficial y se reúne semanalmente. &
Las revisiones de cada incremento pueden programarse con el administrador; para validar con ventanilla y rentas, él debe convocar a esos usuarios. &
La priorización frecuente se complica: el administrador atiende una vez por semana y los usuarios finales no pueden participar libremente en cada ciclo. \\ \hline
K3. Tipo de sistema y criticidad &
Al manejar cobros y datos ciudadanos, una especificación formal ayuda a documentar reglas y controles antes de construir el sistema web. &
Permite probar por partes los trámites, las tasas y el acceso a datos, detectando errores antes de la instalación prevista para el mes 6. &
Las entregas cortas permitirían mostrar pantallas, pero no bastaría una lista de tareas sin pruebas y registros formales sobre cobros y datos. \\ \hline
K4. Riesgo principal &
La base catastral antigua no tiene documentación. Si su integración se analiza tarde, puede comprometer la entrega del sistema en el mes 6. &
Un incremento técnico temprano puede comprobar cómo leer los datos catastrales y reducir el riesgo de descubrir incompatibilidades al final. &
Un prototipo temprano también ayudaría, aunque las decisiones sobre acceso y cambios deben pasar por el canal contractual del GAD. \\ \hline
K5. Documentación y contrato &
Se ajusta a los hitos pactados: ERS aprobada en el mes 2, sistema probado en el mes 6 y acta de entrega en el mes 7. &
Puede mantener la ERS y la trazabilidad actualizadas en cada ciclo y conservar las aprobaciones exigidas en los pliegos. &
El backlog resulta útil para organizar tareas, pero no reemplaza la ERS aprobada ni las actas y pruebas exigidas para cobrar los hitos. \\ \hline
\end{tabularx}
\end{center}

\subsection{Recomendación y riesgos del escenario 3}

\begin{center}
\small
\begin{tabularx}{\textwidth}{|>{\raggedright\arraybackslash}p{4.5cm}|L|}
\hline
\rowcolor{verdeuteq}
\textcolor{white}{\textbf{Elemento}} & \textcolor{white}{\textbf{Respuesta}} \\ \hline
Modelo recomendado (puro o híbrido; si es híbrido, qué parte sigue cada modelo) &
Se propone un enfoque híbrido: control secuencial para aprobar la ERS, gestionar cambios del contrato y cerrar los hitos; trabajo iterativo e incremental para comprobar la integración catastral y desarrollar los módulos de trámites y cobros. \\ \hline
Razón 1 (dato del escenario, criterio y consecuencia) &
K5: el contrato exige una ERS aprobada en el mes 2 y tres entregables con pago asociado. Por ello, conviene una línea base formal que permita demostrar qué se acordó en cada hito. \\ \hline
Razón 2 &
K1: la jefa de rentas ya cambió dos veces el flujo de patentes y se prevé una posible ordenanza en el mes 4. Trabajar por ciclos ayuda a revisar esas partes sin rehacer todo el sistema. \\ \hline
Razón 3 &
K4: la base catastral antigua carece de documentación. Probar la integración en un incremento temprano reduce la posibilidad de descubrir un bloqueo técnico cerca de la entrega del mes 6. \\ \hline
Por qué se descartan los otros dos modelos &
No se recomienda cascada pura: congelar todo en el mes 2 dejaría poco margen ante los cambios de patentes, tasas e integración. Tampoco se recomienda ágil puro: el cliente no participa a diario y el contrato exige documentos, aprobaciones y trámites formales para cambiar el alcance. \\ \hline
Tensión principal del escenario (el dato que apunta a otro modelo) &
El contrato empuja hacia la planificación secuencial, pero las modificaciones del flujo de patentes, la posible ordenanza y la base catastral desconocida requieren revisión iterativa. \\ \hline
Riesgo 1 del modelo elegido y su mitigación &
Riesgo: que cada incremento introduzca cambios que contradigan la ERS aprobada. Mitigación: mantener una matriz de trazabilidad y registrar con el administrador cada solicitud de cambio antes de modificar el alcance contractual. \\ \hline
Riesgo 2 del modelo elegido y su mitigación &
Riesgo: que la integración con el catastro retrase las pruebas del mes 6. Mitigación: hacer una prueba técnica de lectura de datos al comienzo, documentar los campos disponibles y reservar tiempo para corregir incompatibilidades. \\ \hline
\end{tabularx}
\end{center}
